\documentclass[11pt]{article}
\usepackage[a4paper,margin=1in]{geometry}
\usepackage{fontspec}
\usepackage[boldfont=false]{xeCJK}
\setCJKmainfont{FandolSong-Regular.otf}
\usepackage{amsmath}
\usepackage{amssymb}
\usepackage{array}
\usepackage{booktabs}
\usepackage{graphicx}
\usepackage{adjustbox}
\usepackage{enumitem}
\setlist[itemize]{topsep=2pt,partopsep=0pt,parsep=0pt,itemsep=2pt,leftmargin=1.4em}
\setlist[enumerate]{topsep=2pt,partopsep=0pt,parsep=0pt,itemsep=2pt,leftmargin=1.6em}
\usepackage{parskip}
\usepackage{fvextra}
\fvset{breaklines=true,breakanywhere=true,fontsize=\small}
\setlength{\parindent}{0pt}

\begin{document}

\begin{center}
  {\LARGE\bfseries Nuclear physics}\\[0.35em]
  {\large Igcse Physics}
\end{center}
\vspace{0.6em}

\section*{The atom}

An \textbf{atom} 原子 is made of a tiny central \textbf{nucleus} 原子核 with \textbf{electrons} 电子 moving around it (in \textbf{orbit} 轨道, like planets around the Sun).

\begin{itemize}
\item The nucleus has a positive \textbf{charge} 电荷.

\item The electrons have a negative charge.

\item The atom as a whole is neutral, because the positive and negative charges are equal.

\end{itemize}

Almost all the \textbf{mass} 质量 is in the nucleus, but the nucleus is very small compared with the whole atom. So an atom is mostly empty space.

\subsection*{Ions}

An atom is neutral, but it can gain or lose electrons to become an \textbf{ion} 离子.

\begin{itemize}
\item Lose one or more electrons → a positive ion (now there are more protons than electrons).

\item Gain one or more electrons → a negative ion.

\end{itemize}

\subsection*{The alpha-scattering experiment}

This experiment gave the evidence for the nuclear model. \textbf{Alpha particles} α粒子 (small, fast, positive) were fired at a very thin \textbf{gold foil} 金箔, and the \textbf{scattering} 散射 (the way they bounced off) was watched.

The results and what they tell us:

\begin{itemize}
\item Almost all the alpha particles went straight through. → The atom is mostly empty space.

\item A few were \textbf{deflected} 偏转 (bent) through small angles. → The nucleus has a positive charge, which pushes the positive alpha particles away.

\item A very few bounced almost straight back. → The nucleus is very small and very heavy, and holds most of the mass of the atom.

\end{itemize}

\section*{Inside the nucleus}

The nucleus is made of two kinds of particle, together called \textbf{nucleons} 核子:

\begin{itemize}
\item \textbf{protons} 质子, which have a \textbf{relative} 相对 charge of $+1$;

\item \textbf{neutrons} 中子, which have a relative charge of $0$ (they are neutral).

\end{itemize}

An electron has a relative charge of $-1$. A proton and a neutron each have a \textbf{relative mass} of about $1$; an electron is almost massless in comparison.

Two numbers describe a nucleus:

\begin{itemize}
\item the \textbf{proton number} 质子数 $Z$ (also called the atomic number) — the number of protons;

\item the \textbf{nucleon number} 核子数 $A$ (also called the mass number) — the number of protons plus neutrons.

\end{itemize}

So the number of neutrons is $A - Z$.

We write a nucleus in \textbf{nuclide} 核素 notation:

\[
^{A}_{Z}\text{X}
\]

where X is the chemical symbol. For example, $^{197}_{\ 79}\text{Au}$ has $79$ protons and $197 - 79 = 118$ neutrons. The relative charge of the whole nucleus is just $+Z$ (here $+79$), and its relative mass is about $A$.

\subsection*{Isotopes}

\textbf{Isotopes} 同位素 are atoms of the same element (the same $Z$) but with different numbers of neutrons (different $A$). They behave the same in chemistry but differently in the nucleus. For example, $^{12}_{\ 6}\text{C}$ and $^{14}_{\ 6}\text{C}$ are both carbon.

\subsection*{Nuclear fission and fusion}

In \textbf{nuclear fission} 核裂变, a heavy nucleus absorbs a neutron and then splits into two smaller nuclei, giving out two or three neutrons and a lot of \textbf{energy} 能量:

\[
^{235}_{\ 92}\text{U} + ^{1}_{0}\text{n} \rightarrow\ ^{141}_{\ 56}\text{Ba} + ^{92}_{36}\text{Kr} + 3\,^{1}_{0}\text{n}
\]

The top numbers (nucleon numbers) balance on both sides, and so do the bottom numbers (proton numbers). You can use this to find a missing number — for example, how many neutrons are released.

In \textbf{nuclear fusion} 核聚变, two light nuclei join to make a heavier one, also giving out energy. This is how the Sun makes its energy, joining \textbf{hydrogen} 氢 nuclei to make \textbf{helium} 氦:

\[
^{2}_{1}\text{H} + ^{3}_{1}\text{H} \rightarrow\ ^{4}_{2}\text{He} + ^{1}_{0}\text{n}
\]

In both fission and fusion, a small amount of mass is lost and turned into energy.

\section*{Radioactivity}

A nucleus that is \textbf{unstable} 不稳定 will sooner or later break down and give out \textbf{radiation} 辐射. This is \textbf{radioactive} 放射性 \textbf{decay} 衰变. A nucleus may be unstable because it has too many neutrons, or because it is too heavy.

Decay is \textbf{spontaneous} 自发 (it happens on its own, and you cannot speed it up or slow it down) and \textbf{random} 随机 (you cannot say which nucleus will decay next, or exactly when).

\subsection*{Background radiation}

Some radiation is around us all the time. This is \textbf{background radiation} 背景辐射. Its main sources are:

\begin{itemize}
\item \textbf{radon} 氡 gas in the air (usually the biggest source);

\item rocks and buildings;

\item food and drink;

\item \textbf{cosmic rays} 宇宙射线 from space.

\end{itemize}

\subsection*{Measuring radiation}

Radiation can be measured with a \textbf{detector} 探测器 joined to a \textbf{counter} 计数器. The \textbf{count rate} 计数率 is the number of counts each second (or each minute).

To find the true count rate from a source, first measure the background count rate on its own, then subtract it. The answer is the \textbf{corrected} 修正 count rate:

\[
\text{corrected count rate} = \text{measured count rate} - \text{background count rate}
\]

\section*{The three types of radiation}

The radiation can be one of three types. Each type is \textbf{ionising} 电离, which means it can knock electrons off atoms in its path.

\begin{center}
\adjustbox{max width=\linewidth}{%
\begin{tabular}{llll}
\toprule
\textbf{Type} & \textbf{What it is} & \textbf{Ionising effect} & \textbf{Stopped by} \\
\midrule
alpha (α) & a helium nucleus: 2 protons + 2 neutrons, charge $+2$ & strongest & a sheet of paper, or a few cm of air \\
beta (β) & a fast-moving electron, charge $-1$ & medium & a few mm of \textbf{aluminium} 铝 \\
gamma (γ) & a high-energy \textbf{electromagnetic wave} 电磁波, no charge & weakest & thick \textbf{lead} 铅 or concrete (only reduced, never fully stopped) \\
\bottomrule
\end{tabular}}
\end{center}

So an \textbf{alpha particle} is the most \textbf{ionising} but the least \textbf{penetrating} 穿透 (it is easily \textbf{absorbed} 吸收). A \textbf{beta particle} β粒子 is in the middle. \textbf{Gamma radiation} γ射线 is the least ionising but the most penetrating.

We can explain the ionising effects from charge and \textbf{kinetic energy} 动能: an alpha particle has a large charge ($+2$) and is slow and heavy, so it pulls strongly on the electrons it passes and ionises a lot. A beta particle has a smaller charge and moves faster, so it ionises less.

\subsection*{Deflection in fields}

Because alpha and beta particles are charged, they are deflected by an \textbf{electric field} 电场 and by a \textbf{magnetic field} 磁场. They bend in \emph{opposite} directions, because their charges have opposite signs, and the lighter beta particle bends more. Gamma rays have no charge, so they are not deflected at all.

\section*{Decay equations}

When a nucleus decays, the nucleon and proton numbers must balance on both sides.

\textbf{Alpha decay} — the nucleus loses 2 protons and 2 neutrons, so $A$ falls by $4$ and $Z$ falls by $2$. It becomes a \emph{different element}:

\[
^{A}_{Z}\text{X} \rightarrow\ ^{A-4}_{Z-2}\text{Y} + ^{4}_{2}\alpha
\]

\textbf{Beta decay} — inside the nucleus a neutron changes into a proton plus an electron:

\[
\text{neutron} \rightarrow \text{proton} + \text{electron}
\]

The fast electron leaves as the beta particle. So $A$ stays the same but $Z$ rises by $1$, giving a different element:

\[
^{24}_{11}\text{Na} \rightarrow\ ^{24}_{12}\text{Mg} + ^{\ \ 0}_{-1}\beta
\]

\textbf{Gamma emission} — the nucleus loses only energy, so $A$ and $Z$ do not change. Alpha and beta decay leave the nucleus more \textbf{stable} 稳定; gamma is often given out at the same time to carry away spare energy.

\section*{Half-life}

Because decay is random, we cannot follow one nucleus. Instead we describe a large \textbf{sample} 样品 using its half-life.

The \textbf{half-life} 半衰期 of an isotope is the time taken for half the unstable nuclei in a sample to decay. After each half-life, the count rate (or the number of unstable nuclei left) falls to half.

For example, if a source has a count rate of $800$ counts/min and a half-life of $3$ hours:

\begin{center}
\adjustbox{max width=\linewidth}{%
\begin{tabular}{lllll}
\toprule
\textbf{Time / hours} & \textbf{0} & \textbf{3} & \textbf{6} & \textbf{9} \\
\midrule
Count rate / (counts/min) & 800 & 400 & 200 & 100 \\
\bottomrule
\end{tabular}}
\end{center}

After $6$ hours (two half-lives) the count rate has halved twice: $800 \to 400 \to 200$. You can read a half-life off a decay graph by finding the time for the count rate to drop from any value to half of it.

\section*{Uses of radioactivity}

The isotope chosen for a job depends on its type of radiation and its half-life:

\begin{itemize}
\item \textbf{smoke alarms} 烟雾报警器 use an alpha source (alpha is easily blocked, so it is safe outside the alarm, and a long half-life means it lasts for years);

\item using radiation on food, or on equipment, kills \textbf{bacteria} 细菌 and germs — this needs penetrating gamma rays to \textbf{sterilise} 消毒 the sealed item;

\item measuring and controlling the thickness of paper or metal sheets uses a beta source (the amount that passes through changes with the thickness);

\item gamma rays are used to find and to treat \textbf{cancer} 癌症 inside the body, because they can pass out of (or into) the body.

\end{itemize}

\section*{Safety}

Ionising radiation harms \textbf{living cells} 细胞. It can cause cell death, \textbf{mutations} 突变 (changes to the genes) and cancer.

When working with radioactive sources, keep the \textbf{dose} 剂量 (the amount of radiation received) low by:

\begin{itemize}
\item reducing the \textbf{exposure} 照射 time — spend as little time near the source as possible;

\item increasing the distance between you and the source;

\item using \textbf{shielding} 屏蔽, such as lead, between the source and your body.

\end{itemize}

Radioactive sources should be handled with tongs (not bare hands), kept pointing away from people, and stored in a lead-lined box.


\end{document}
